%\section{Introduction}
Speech representation learning has seen remarkable progress with the advent of self-supervised learning (SSL) techniques, enabling the development of foundation models capable of understanding and processing human speech. While state-of-the-art (SOTA) speech models have achieved unprecedented success in high-resource languages, low-resource languages like Bengali often suffer from a lack of dedicated, efficient models built from scratch. Most existing works in Bengali speech AI rely on fine-tuning massive, pre-trained multilingual models, which are computationally expensive and not optimized specifically for the linguistic nuances of Bengali. Furthermore, the immense size of these SOTA models limits their deployability in resource-constrained environments. This research addresses this gap by developing an AI/ML speech foundation model specifically for Bengali, trained entirely from scratch using self-supervised learning on publicly available datasets. By doing so, we demonstrate that training a compact, dedicated model from scratch is not only viable but can also yield highly effective representations for Bengali speech.

\section{Problem Statement}
Despite the growing interest in speech technologies, there is a significant scarcity of foundational speech models built from scratch for the Bengali language. Current approaches predominantly involve fine-tuning extremely large SOTA models trained on other languages or multilingual corpora. This reliance on massive pre-trained models presents two major challenges: first, the computational cost and memory footprint of these models make them impractical for widespread deployment or research with limited resources; second, they often lack specialized adaptation to the unique phonetic and phonological characteristics of Bengali. There are currently no researchers who have successfully developed a small-scale, highly efficient speech foundation model for Bengali from scratch using self-supervised learning on public datasets. Therefore, there is a critical need to investigate whether a significantly smaller, dedicated model can be trained from scratch to achieve competitive or superior performance in specific metrics compared to existing heavy-weight SOTA models.

\section{Thesis Objective}
The primary objective of this thesis is to design, train, and evaluate a Continuous Latent Encoder (CLE) for Bengali speech representation learning. By developing a compact speech foundation model from scratch using self-supervised learning techniques, we aim to provide a highly efficient alternative to fine-tuning large SOTA models. The specific objectives of this research are summarized as follows:

\begin{itemize}
    \item To develop a lightweight speech foundation model (a 23.8M-parameter continuous-latent autoencoder) specifically tailored for 16 kHz Bengali speech using publicly available datasets.
    \item To demonstrate that training a model from scratch using self-supervised learning is a viable and highly effective pathway for Bengali speech research, eliminating the strict dependency on fine-tuning massive multilingual models.
    \item To evaluate the performance of the proposed compact model and show that despite its significantly small size compared to SOTA speech models, it achieves superior results in specific evaluation metrics.
    \item To provide a robust, frame-level speech representation that can be utilized for various downstream Bengali speech analysis tasks, establishing a foundation for future research in low-resource language speech AI.
\end{itemize}

\section{Thesis Outline}
The rest of the paper has been organized as follows: Chapter-\hyperref[related_work]{2} explains some relevant research in speech foundation models and self-supervised learning, along with their contributions and limitations. Chapter-\hyperref[methodology]{3} describes the methodology, dataset, and the continuous-latent autoencoder architecture used in this research in detail. Chapter-\hyperref[result_and_analysis]{4} presents the results we obtained from this study, including comparative analyses with SOTA models to evaluate the effectiveness of our approach. We also describe some limitations and future goals of this paper. Finally, in Chapter-\hyperref[conclusion]{5}, we conclude the thesis, providing the study's findings and potential areas for further exploration.

% \nomenclature{$ODI$}{One day International}
% \nomenclature{$T20$}{Twenty Twenty}
% \nomenclature{$ODI$}{One day International}
% \nomenclature{$IPL$}{Indian Premier League}
% \nomenclature{$ODI$}{One day International}
% \nomenclature{$MR$}{Runs scored by Home team}
% \nomenclature{$OR$}{Runs scored by the opponent team}
% \nomenclature{$MRN$}{Home Team Run Rate}
% \nomenclature{$ORN$}{Opponent Team Run Rate}
% \nomenclature{$LBW$}{Leg before Wicket}